\chapter*{绪论：智能状态动力学的研究对象}
\addcontentsline{toc}{chapter}{绪论：智能状态动力学的研究对象}
\markboth{绪论}{智能状态动力学的研究对象}
\label{intro:scope}

\section*{问题与理论层级}
我们研究智能系统在有限资源下的持续交互过程。系统论使我们区分整体、部件和环境；协同论提示我们寻找能够组织局部活动的低维变量；耗散结构理论提示我们注意开放交换与持续维持；复杂系统理论提供非线性、网络、分岔与跨尺度视角；认知科学约束记忆、注意、控制和自我模型的具体解释。上述来源提供研究原则，而不是直接给出唯一的智能方程。

HSF-HD 承担中间层的工作：明确状态空间，提出更新机制，推导有条件的性质，建立可观测指标。AgentOS 承担实例化工作：将状态、历史、结构、控制和资源约束落实为可运行的记忆与调度系统。不同架构可以实现同一个动力学模型族，同一个架构也可以承载不同模型。

\section*{最小模型族与对象类型}
令 $\mathbf{x}\in\mathbb{R}^{d_x}$ 表示快速活动状态，$\mathbf{m}\in\mathbb{R}^{d_m}$ 表示历史及中间记忆，$S\in\mathcal{S}$ 表示持久结构，$\mathbf{v}\in\mathbb{R}^{d_v}$ 表示目标与偏好状态。若结构是矩阵，则 $\mathcal{S}$ 可以是约束矩阵集合；若结构是图，则其离散改变由事件更新定义。下面的微分形式只用于结构具有连续参数化的情形：
\begin{align}
\dot{\mathbf{x}}&=F(\mathbf{x},\mathbf{m},S,\mathbf{v},\mathbf{u},\mathbf{c}),\\
\dot{\mathbf{m}}&=\varepsilon_m H(\mathbf{x},\mathbf{m},S,\mathbf{v}),\\
\dot{S}&=\varepsilon_S G(\mathbf{x},\mathbf{m},S,\mathbf{v}),\\
\dot{\mathbf{v}}&=\varepsilon_v Q(\mathbf{x},\mathbf{m},S,\mathbf{v},\mathbf{u}),\\
\mathbf{a}&=R(\mathbf{x},\mathbf{m},S,\mathbf{v}).
\end{align}
这些等式定义一个接口，而不是宣称所有智能系统具有相同的右端函数。向量场在工作域内局部 Lipschitz，可提供局部解的存在唯一性；若需要全局解，还须控制增长或证明状态留在有界不变域中。$\varepsilon_m,\varepsilon_S,\varepsilon_v$ 指定相对于快速状态的更新速率，只有满足尺度分离条件时才可采用准静态近似。

输入 $\mathbf{u}$ 必须经过编码进入状态空间，控制 $\mathbf{c}$ 由观测和策略生成。输出动作还需通过环境转移与感知返回，形成闭环。观测不完整、通信延迟和计算预算都可能改变闭环性质。

\section*{从闭环到验证}
给定模型后，我们首先检查维度和定义域，再研究有界性、平衡点、收敛和扰动响应。随后将模型离散化，测量实现中的活动强度、记忆误差、更新代价、任务表现与干预效果。任务成功不由低状态范数自动推出，内部预测改善也不由输出流畅自动推出。

三卷依次讨论表示、演化与工程。第一卷的形质分解建立对象；第二卷的传播、学习与控制建立过程；第三卷的比较与实现检查这些机制在不同载体上的成立程度。我们保留数学工具的多样性，但用统一的对象类型和可验证条件约束它们。

\section*{为什么状态动力学需要独立的理论层}
我们从一个具体的计算过程出发。一个智能体需要连续阅读一组文档，发现相互矛盾的陈述，保留证据，并在资源预算内形成报告。单次输出质量不能描述这个过程的全部性质。相同的编码器在不同历史、工作记忆、工具结果和控制策略下，会形成不同的行动路径。我们需要解释这些路径如何发生、何时可以维持、又在什么条件下偏离目标。

在这个例子中，当前激活的证据及候选结论属于快速状态；已经建立的文档索引与概念关系属于持久结构；先前读取的位置、待核实事项和中间结论属于记忆；任务的优先级、容许误差及预算偏好属于目标与价值状态。一个文档的文本内容与其来源关系需要共同保存，否则系统可能记住陈述而遗失证据身份。这个问题连接了形质绑定、状态传播、历史保持和运行时调度，不能仅用某一个模块的输入输出来解释。

因此，我们把 HSF-HD 组织为模型族。它为不同实现提供共同的对象与约束，而不预先规定所有智能体采用同一种神经网络、同一个维度或同一种时间参数。我们的“空气动力学”类比强调这种一般性：研究对象是运行中的状态及其与结构、环境和控制的耦合。类比本身不赋予模型任何物理定律。

\section*{由连续模型到实际更新}
令联合状态 $z=(\mathbf{x},\mathbf{m},S,\mathbf{v})$，并把各分量的更新写为 $\dot z=f(z,\mathbf{u},\mathbf{c})$。这里要求各分量已经在同一局部坐标系统中定义；图结构的增删等离散变化仍通过单独的事件规则处理。对于固定输入、固定控制和固定步长 $h>0$，显式 Euler 近似为
\begin{equation}
 z_{k+1}=z_k+h f(z_k,\mathbf{u}_k,\mathbf{c}_k).
\end{equation}
这是一种数值实现选择，而不是从连续形式直接得到的稳定性结论。若 $f$ 在所考察区域对 $z$ 是 $L$-Lipschitz，并且两条轨迹使用同一输入与控制，则一步扰动满足
\begin{align}
 \|z_{k+1}-\widetilde z_{k+1}\|
 &\leq \|z_k-\widetilde z_k\|
       +h\|f(z_k,\mathbf{u}_k,\mathbf{c}_k)
           -f(\widetilde z_k,\mathbf{u}_k,\mathbf{c}_k)\|\\
 &\leq (1+hL)\|z_k-\widetilde z_k\|.
\end{align}
这个上界描述误差放大的可能范围，不能推出收缩。若两个智能体由于状态差异而选择不同动作，还需要把策略与环境反馈一起纳入联合系统；不能继续使用“相同控制”的假设。

为看清步长的作用，我们考虑 $\dot{\mathbf{x}}=-A\mathbf{x}$，其中 $A=A^\top\succ0$。连续解在各特征方向上按 $e^{-\lambda_i t}$ 衰减。离散更新为 $\mathbf{x}_{k+1}=(I-hA)\mathbf{x}_k$；其渐近稳定条件是所有特征因子满足 $|1-h\lambda_i|<1$，从而
\begin{equation}
 0<h<\frac{2}{\lambda_{\max}(A)}.
\end{equation}
我们由此区分模型的连续性质与实现的离散性质。异步调用、延迟、噪声及变化的结构会改变这个分析，需要在对应章节中明确加入。

\section*{结构、活动与历史如何共同工作}
对有限关系图，我们用 $S$ 表示邻接或耦合参数，用 $X$ 表示节点上的活动。若 $S$ 是对称非负权图，令 $L_S=D_S-S$，其中 $D_S$ 的对角元是各节点的权重和。对任意向量 $q$，我们直接计算
\begin{equation}
 q^\top L_S q=\frac12\sum_{i,j}S_{ij}(q_i-q_j)^2\geq0.
\end{equation}
所以 $-L_S X$ 可以描述关系结构上的差异消减。然而，差异消减既不是任务正确性，也不等于语义理解。若所有节点的错误状态趋于一致，系统仍然可能失败。输入项、目标项、约束和读出必须与传播共同定义。

我们可以选择一个具体的快速过程
\begin{equation}
 \dot X=-L_S X-\gamma X+B\mathbf{u}+U\mathbf{c},\qquad \gamma>0,
\end{equation}
并令历史更新为 $\dot{\mathbf{m}}=-\lambda\mathbf{m}+R(X)$，其中 $\lambda>0$。对给定的 $X(t)$，历史方程的解是
\begin{equation}
 \mathbf{m}(t)=e^{-\lambda t}\mathbf{m}(0)
 +\int_0^t e^{-\lambda(t-s)}R(X(s))\,ds.
\end{equation}
这说明历史变量可以实现带衰减权重的累积。它保存哪些信息取决于 $R$，保存多久取决于 $\lambda$；一个积分表达式本身不保证事件顺序、来源身份或完整情节已经被保存。对于文档任务，我们还需要显式的来源键、时间索引与外部证据记录。

慢速结构更新可以由目标函数驱动，例如 $\dot S=-\varepsilon_S\nabla_S\mathcal J(X,S,\mathbf{m},\mathbf{v})$。但只有在其他变量固定、约束处理正确时，结构方向上的导数才给出 $-\varepsilon_S\|\nabla_S\mathcal J\|^2$。若整个系统同时更新，则
\begin{equation}
 \frac{d\mathcal J}{dt}
 =\langle\nabla_X\mathcal J,\dot X\rangle
 +\langle\nabla_S\mathcal J,\dot S\rangle
 +\langle\nabla_{\mathbf{m}}\mathcal J,\dot{\mathbf{m}}\rangle
 +\langle\nabla_{\mathbf{v}}\mathcal J,\dot{\mathbf{v}}\rangle
 +\partial_t\mathcal J.
\end{equation}
因此，我们不能仅凭某一分量采用梯度更新，就宣布整个耦合系统的目标单调下降。这一原则贯穿本书的跨尺度分析。

\section*{从数学性质到任务证据}
我们分别考察四个层面的结果。首先，更新规则是否定义良好，变量维度、输入范围与边界条件是否一致。其次，轨迹是否有界、是否存在平衡状态、在什么扰动下能够恢复。再次，保存的表示是否支持来源追溯、目标预测和行动选择。最后，系统是否在独立任务与资源预算下达到预期表现。

以文档任务为例，状态范数有界只能说明活动没有在模型中无限增长。它不能保证报告中的结论被证据支持。我们还需要检查引文的来源一致性、矛盾处理的正确率、遗失证据的比例、预算超限的次数以及任务中断后的恢复能力。每个观测指标应对应一种明确机制，避免把一个好看的内部曲线解释为整个智能已经成立。

我们提出的机制也需要比较对象。历史衰减应与无历史及不同衰减长度比较；结构更新应与冻结结构比较；控制应与固定策略比较。若增加控制后结果改善，却同时增加了大量计算预算，就需要进一步区分机制效果与资源效果。这里的比较原则是理论验证的一部分。

\section*{HSF-HD 与 AgentOS 的接口}
我们把基础理论、状态动力学与运行时工程之间的关系写为三个层次：基础系统科学与认知科学提出组织原则及经验问题；HSF-HD 定义状态、结构、记忆、目标及其耦合模型；AgentOS 在具体计算平台上落实这些模型的存储、调度、控制和交互。

这个关系允许双向修正。AgentOS 的观测可以暴露某个动力学假设的不足，例如工具延迟破坏同步近似，记忆压缩破坏来源绑定，或者结构更新引起原有稳定性边界变化。我们据此修正 HSF-HD 的模型条件，而不是要求工程系统服从一条缺少适用性检验的总方程。

运行时中的对象也不是理论符号的简单改名。一个状态向量可以分布在多个组件中；一项历史变量可以由数据库与摘要共同实现；一次理论上的连续更新可以由多个异步事件近似完成。我们需要给出从实现状态到理论状态的读出映射，以及该映射在哪些时刻有效。自我模型是系统所维护的一种对象，不与全部目标变量或全部运行状态等同。

\section*{全书的推进路径}
第一卷从结构与属性开始，讨论绑定、表示空间、跨域映射及联合表示架构，为后续更新规则提供对象基础。第二卷将这些对象放入交互过程，分析状态传播、控制、记忆、目标和多尺度耦合，逐步区分可以推导的性质与有待检验的机制。第三卷讨论不同系统的组织方式及工程实现，关心理论变量如何被观测、维护和执行。附录保留形式化工具、实现补充和进一步阅读的入口。

我们在各章先展开问题与例证，再建立变量与假设，随后给出机制和推导，最后讨论边界与实现后果。丰富的论述不是与公式无关的装饰：它使读者理解为什么引入这些对象，以及一个数学结果究竟回答了哪个问题。相应地，公式也不能只作为论述的象征；每一个运算都必须有明确对象，每一个结论都必须有足够条件。
